\chapter{Supraventricular Tachycardia}
\index{Supraventricular Tachycardia}
\label{sec:Supraventricular Tachycardia}

\section{Overview}
Supraventricular tachycardia (SVT) is a paroxysmal arrhythmia characterized by a sudden onset of rapid heart rate originating above the ventricles, typically between 150--250 beats per minute. It results from abnormal electrical pathways or reentrant circuits involving the atrioventricular (AV) node, atrial tissue, or accessory conduction pathways. While rarely life-threatening, SVT episodes can cause significant discomfort, anxiety, and impairment of daily function.
\section{Classification}

\begin{description}[font=\normalfont\bfseries,leftmargin=3em,labelwidth=2.5em]
  \item[Cause] Idiopathic or structural
  \item[Organ System] Heart
  \item[Pathophysiology] Reentrant circuit or abnormal automaticity in atrial tissue or AV node producing paroxysmal tachycardia
  \item[Duration] Acute (paroxysmal episodes lasting seconds to hours)
  \item[Onset] Sudden
  \item[Mode of Transmission] Non-communicable
  \item[Clinical Course] Episodic with abrupt onset and termination
  \item[Severity] Mild to Moderate
  \item[Extent of Disease] Localized (heart conduction system)
  \item[Age Group] All ages; more common in young adults and women
  \item[Epidemiology] Prevalence approximately 2.25 per 1000 in the general population
  \item[ICD-11 Code] BC94.1
\end{description}

\section{Causes}
SVT is most commonly caused by a reentrant tachycardia involving dual AV nodal pathways (AV nodal reentrant tachycardia, AVNRT) accounting for approximately 60\% of cases. Accessory pathways such as in Wolff-Parkinson-White syndrome create an additional conduction route that supports reentrant circuits (AV reentrant tachycardia, AVRT). Atrial tachycardia, focal or multifocal, arises from abnormal automaticity or micro-reentry within atrial tissue.

\section{Risk Factors}

\begin{itemize}[noitemsep]
  \item Female sex (higher incidence of AVNRT)
  \item Preexcitation syndromes such as Wolff-Parkinson-White (WPW) pattern
  \item Anxiety, stress, or panic attacks
  \item Caffeine, alcohol, or stimulant use
  \item Structural heart disease (e.g., congenital heart disease, valvular disease)
\end{itemize}

\section{Signs and Symptoms}

\subsection{Early symptoms}
\begin{itemize}[noitemsep]
  \item Early manifestations of supraventricular tachycardia may be subtle and nonspecific
\end{itemize}

\subsection{Common symptoms}
\begin{itemize}[noitemsep]
  \item \subsection{Common symptoms}
  \item \textbf{Common symptoms:}
  \item Palpitations or racing heartbeat with abrupt onset and termination
  \item Lightheadedness or dizziness during episodes
  \item Shortness of breath
  \item Chest fluttering or discomfort
  \item Fatigue or weakness following episodes
  \item Near-syncope or syncope in some cases
  \item Anxiety or sense of impending doom
  \item Polyuria (due to atrial natriuretic peptide release)
  \item Pain or discomfort in the affected area
  \item Difficulty performing daily activities
\end{itemize}

\subsection{Less common symptoms}
\begin{itemize}[noitemsep]
  \item Atypical or less common presentations of supraventricular tachycardia may occur
\end{itemize}

\subsection{Emergency warning signs}
\begin{itemize}[noitemsep]
  \item Severe or worsening symptoms of supraventricular tachycardia require immediate medical evaluation
\end{itemize}
\section{Progression and Stages}
Episodes of SVT typically begin suddenly and can last from a few seconds to several hours before terminating spontaneously or with intervention. Frequency and duration of episodes vary widely among individuals, with some experiencing rare isolated events and others having frequent daily episodes that significantly affect quality of life. Over time, very frequent or prolonged episodes can sometimes lead to tachycardia-induced cardiomyopathy with reduced left ventricular function.

\section{Complications}

\begin{itemize}[noitemsep]
  \item Tachycardia-induced cardiomyopathy from frequent or sustained episodes
  \item Syncope or near-syncope during rapid episodes
  \item Preexcited atrial fibrillation in WPW patients with risk of degeneration to ventricular fibrillation
  \item Reduced quality of life due to recurrent symptoms
  \item Emotional and psychological effects
\end{itemize}

\section{Diagnosis}

\begin{itemize}[noitemsep]
  \item Electrocardiogram (ECG) during an episode showing narrow-complex tachycardia
  \item Holter monitor or event recorder for episodic capture
  \item Electrophysiology study for definitive diagnosis and ablation planning
  \item Echocardiogram to evaluate for structural heart disease
  \item Lab tests to rule out thyroid disorders and electrolyte abnormalities
  \item Imaging tests (rarely needed beyond echocardiogram)
\end{itemize}

\section{Prevention}

\begin{itemize}[noitemsep]
  \item Avoid known triggers such as caffeine, alcohol, and stimulants
  \item Manage stress through relaxation and mindfulness techniques
  \item Maintain adequate hydration and electrolyte balance
  \item Catheter ablation as curative therapy for recurrent SVT
  \item Vagal maneuvers taught for acute episode termination
  \item Go for regular check-ups
  \item Follow your doctor's screening recommendations
\end{itemize}

\section{Treatment Options}

\begin{itemize}[noitemsep]
  \item Vagal maneuvers (Valsalva, carotid sinus massage) for acute termination
  \item Adenosine intravenous bolus for acute episode conversion
  \item Calcium channel blockers or beta-blockers for rate control and prevention
  \item Catheter ablation (radiofrequency or cryoablation) for definitive treatment
  \item Lifestyle changes and self-care
  \item Surgery when needed
\end{itemize}

\section{Living with It}

\begin{itemize}[noitemsep]
  \item Follow your treatment plan as prescribed
  \item Keep up with regular appointments
  \item Adjust your daily habits as needed
  \item Reach out to family, friends, or support groups
  \item Keep learning about your condition
\end{itemize}

\section{Outlook}
The prognosis for SVT is excellent, as it is rarely life-threatening and can often be cured definitively with catheter ablation (success rate exceeding 95\% in most cases). Patients with managed or treated SVT can expect a normal lifespan and quality of life with minimal restrictions.
